\section{Założenia projektowe – wymagania}		%1
\begin{enumerate}
\item Założyć domenę \texttt{NIS}, skonfigurować odpowiednio serwer dodać użytkowników domeny, ustawić dla nich różny poziom dostępu.
\item Dołączyć hosty do domeny \texttt{NIS}.
\item Uruchomić na serwerze i każdym z hostów serwer \texttt{NFS}.
\item Udostępnić zasoby na różnych prawach dla pozostałych hostów (jeden katalog pełny dostęp dla wszystkich, drugi dostęp ograniczony wg uznania).
\item Zmontować udostępnione katalogi na hostach innych hostach. (katalogi udostępnione na pełnych prawach montowanie automatyczne przy starcie, pozostałe zamontować ręcznie).
\item Udokumentować wszystkie kroki w postaci zrzutów ekranu oraz opisów. 
\end{enumerate}
\begin{itemize}
      \item \textbf{Tworzenie maszyn wirtualnych}
      \begin{itemize}
          \item \textit{Wymagane maszyny:}
          \begin{itemize}
              \item \textbf{1×} Serwer \texttt{NIS}/\texttt{NFS} (główny węzeł)
              \item \textbf{2×} Hosty klienckie (dołączone do domeny \texttt{NIS})
          \end{itemize}
          \item \textit{Wspólne ustawienia dla wszystkich VM:}
          \begin{itemize}
              \item Typ systemu: \texttt{Linux} / \texttt{Ubuntu} (64-bit)
              \item Wersja Ubuntu: \texttt{ubuntu-24.04.2-live-server-amd64}
              \item Tryb instalacji: \texttt{Server} (opcjonalne GUI Gnome)
          \end{itemize}
          \item \textit{Minimalne wymagania sprzętowe:}
          \begin{itemize}
              \item Pamięć RAM: 
              \begin{itemize}
                  \item Serwer: \textbf{4096 MB} 
                  \item Klienty: \textbf{2048 MB}
              \end{itemize}
              \item Procesor: 
              \begin{itemize}
                  \item Serwer: \textbf{2} rdzenie
                  \item Klienty: \textbf{1-2} rdzenie
              \end{itemize}
              \item Dysk twardy (przydzielany dynamicznie): 
              \begin{itemize}
                  \item Serwer: \textbf{20 GB}
                  \item Klienty: \textbf{20 GB}
              \end{itemize}
          \end{itemize}
      \end{itemize}
      \clearpage
      \item \textbf{Konfiguracja sieci}
      \begin{itemize}
          \item \textit{Wymagania:}
          \begin{itemize}
              \item Maszyny muszą komunikować się ze sobą (\texttt{NIS}/\texttt{NFS})
              \item Serwer powinien mieć statyczny IP
          \end{itemize}
          \item \textit{Kroki dla wszystkich VM:}
          \begin{itemize}
              \item Typ karty sieciowej:
              \begin{itemize}
                  \item Karta 1: \texttt{NAT} (dostęp do internetu dla instalacji pakietów)
                  \item Karta 2: Siec wewnętrzna (\texttt{Internal Network}) z nazwą np. \texttt{nis-network}
              \end{itemize}
          \end{itemize}
          \item \textit{Konfiguracja w Ubuntu Server:}
          \begin{itemize}
              \item Podczas instalacji wybierz kartę \texttt{ens3} (\texttt{NAT}) do pobierania pakietów.
              \item Dla karty \texttt{ens7} (\texttt{Internal Network}) ustaw:
              \begin{itemize}
                  \item Serwer: Statyczny IP (np. \textbf{192.168.10.10/24})
                  \item Klienci: \texttt{DHCP} lub statyczne IP (np. \textbf{192.168.10.20}, \textbf{192.168.10.30})
              \end{itemize}
          \end{itemize}
      \end{itemize}
      
      \item \textbf{Dodatkowe ustawienia VirtualBox}
      \begin{itemize}
          \item \textit{Dla serwera:}
          \begin{itemize}
              \item Udostępnianie folderów (opcjonalnie):
              \begin{itemize}
                  \item W \texttt{VirtualBox}: \texttt{Settings → Shared Folders}
                  \item Dodaj folder do współdzielenia np. \texttt{/shared} (przydatne do przenoszenia plików)
              \end{itemize}
              \item Porty przekazywane (jeśli potrzebujesz dostępu z hosta):
              \begin{itemize}
                  \item \texttt{SSH}: \texttt{Settings → Network → Advanced → Port Forwarding}
                  \item Dodaj regułę: Host Port: \textbf{2222} → Guest Port: \textbf{22}
              \end{itemize}
          \end{itemize}
          \item \textit{Dla klientów:}
          \begin{itemize}
              \item Tryb sieciowy:
              \begin{itemize}
                  \item Wyłącz kartę \texttt{NAT} po instalacji pakietów (pozostaw tylko \texttt{Internal Network})
              \end{itemize}
          \end{itemize}
      \end{itemize}
      \end{itemize}